% ============================================================
% 摘要
% ============================================================

% ------------------------------------------------------------
% 中文摘要
% ------------------------------------------------------------
\begin{center}
    {\sanhao\sffamily\bfseries 摘\quad 要}
\end{center}
\addcontentsline{toc}{chapter}{摘要}

\vspace{0.5cm}

随着科学技术的快速发展，学术论文的排版质量越来越受到重视。LaTeX 作为一种专业的排版系统，以其高质量的输出效果和强大的数学公式处理能力，在学术界得到了广泛的应用。本文针对心河大学学位论文的排版需求，设计并实现了一套基于 LaTeX 的论文排版模板系统。

本文首先分析了现有学位论文排版工具存在的问题，阐述了 LaTeX 排版系统的优势。其次，研究了心河大学学位论文的格式规范要求，包括页面设置、字体字号、章节标题、图表公式、参考文献等方面的具体规定。在此基础上，设计了论文模板的整体架构，并实现了文档类文件（.cls）的开发。

本文的主要工作包括：（1）设计了符合心河大学规范的页面布局和字体设置方案；（2）实现了自动编号的章节标题、图表、公式环境；（3）配置了符合GB/T 7714标准的参考文献格式；（4）开发了封面、摘要、目录、致谢等常用页面模板。

通过实际应用测试，本模板能够较好地满足心河大学学位论文的排版需求，生成的PDF文档格式规范、美观大方，具有较高的实用价值。

\vspace{1cm}
\noindent\textbf{关键词：}LaTeX；学位论文；排版模板；GB/T 7714

\newpage

% ------------------------------------------------------------
% 英文摘要
% ------------------------------------------------------------
\begin{center}
    {\sanhao\bfseries Abstract}
\end{center}
\addcontentsline{toc}{chapter}{Abstract}

\vspace{0.5cm}

With the rapid development of science and technology, the quality of academic paper typesetting has received increasing attention. As a professional typesetting system, LaTeX has been widely used in academia due to its high-quality output and powerful mathematical formula processing capabilities. This thesis designs and implements a LaTeX-based thesis template system for the typesetting requirements of domestic university dissertations.

This thesis first analyzes the problems existing in current dissertation typesetting tools and elaborates on the advantages of the LaTeX typesetting system. Secondly, it studies the format specifications of domestic university dissertations, including specific requirements for page settings, fonts, chapter titles, figures, tables, equations, and references. On this basis, the overall architecture of the thesis template is designed, and the document class file (.cls) is developed.

The main contributions of this thesis include: (1) designing a page layout and font setting scheme that conforms to domestic university standards; (2) implementing automatically numbered chapter titles, figures, tables, and equation environments; (3) configuring the reference format according to GB/T 7714 standard; (4) developing templates for common pages such as cover, abstract, table of contents, and acknowledgements.

Through practical application testing, this template can well meet the typesetting requirements of domestic university dissertations. The generated PDF documents are well-formatted and aesthetically pleasing, demonstrating high practical value.

\vspace{1cm}
\noindent\textbf{Keywords:} LaTeX; Dissertation; Typesetting Template; GB/T 7714
